\documentclass[11pt,letterpaper]{article}

\usepackage[utf8]{inputenc}
\usepackage[T1]{fontenc}
\usepackage[spanish,es-nodecimaldot]{babel}
\usepackage[protrusion=true,expansion=false]{microtype}
\usepackage[letterpaper,margin=2.3cm,headheight=15pt]{geometry}
\usepackage{xcolor}
\usepackage{booktabs}
\usepackage{tabularx}
\usepackage{longtable}
\usepackage{enumitem}
\usepackage{fancyhdr}
\usepackage{hyperref}

\ifdefined\pdfinfoomitdate\pdfinfoomitdate=1\fi
\ifdefined\pdftrailerid\pdftrailerid{}\fi
\ifdefined\pdfsuppressptexinfo\pdfsuppressptexinfo=15\fi

\definecolor{ugblue}{HTML}{006A8D}
\definecolor{deepgreen}{HTML}{123D3A}
\definecolor{softgray}{HTML}{F2F5F4}
\definecolor{warning}{HTML}{A45A00}

\hypersetup{
  colorlinks=true,
  linkcolor=ugblue,
  urlcolor=ugblue,
  pdftitle={Bitacora tecnica - Prototipo web de BI asistido por IA},
  pdfauthor={Emily Robles Alvarado y Lesly Velasquez Anchundia}
}

\pagestyle{fancy}
\fancyhf{}
\lhead{Bitacora tecnica}
\rhead{Prototipo BI asistido por IA}
\cfoot{\thepage}

\setlist[itemize]{leftmargin=1.5em,itemsep=0.25em,topsep=0.35em}
\setlength{\parindent}{0pt}
\setlength{\parskip}{0.55em}
\renewcommand{\arraystretch}{1.25}

\newcommand{\statusok}{\textcolor{deepgreen}{\textbf{Completado}}}
\newcommand{\statuspending}{\textcolor{warning}{\textbf{Pendiente}}}
\newcommand{\file}[1]{\texttt{#1}}

\begin{document}
\hypersetup{pageanchor=false}

\begin{titlepage}
  \thispagestyle{empty}
  \vspace*{1.2cm}
  {\color{ugblue}\rule{\textwidth}{2.5pt}}
  \vspace{1.4cm}

  {\Large Documentación técnica del producto\\[0.25em]Prototipo BI asistido por inteligencia artificial}

  \vfill

  {\color{deepgreen}\fontsize{30}{35}\selectfont\textbf{Bitácora técnica}}\\[0.8em]
  {\LARGE Prueba de concepto de un prototipo funcional de BI asistido por IA para la construcción semiautomatizada y supervisada de un datamart de ventas}

  \vspace{1.6cm}

  \begin{tabularx}{\textwidth}{@{}>{\bfseries}l X@{}}
    Responsables & Emily Dayan Robles Alvarado y Lesly Samay Velásquez Anchundia \\
    Tema & Construcción asistida de un datamart, KPI, analítica explicable y reportería \\
    Repositorio local & \file{/home/ceo/Proyectos/BI} \\
    Versión de la bitácora & 0.5.0 \\
    Fecha de apertura & 24 de agosto de 2026 \\
  \end{tabularx}

  \vfill
  {\color{ugblue}\rule{\textwidth}{2.5pt}}
\end{titlepage}

\hypersetup{pageanchor=true}
\pagenumbering{roman}
\tableofcontents
\newpage
\pagenumbering{arabic}

\section{Propósito y reglas de mantenimiento}

Este documento conserva la trazabilidad técnica del prototipo: decisiones, cambios, verificaciones, incidentes, evidencias y trabajo pendiente. Debe actualizarse al final de cada sesión significativa y antes de etiquetar una entrega.

Reglas:

\begin{itemize}
  \item No eliminar incidentes ni resultados fallidos; agregar su resolución cuando exista.
  \item Registrar comandos y resultados relevantes sin copiar contraseñas, tokens o datos sensibles.
  \item Relacionar cada hito con un commit, pull request o etiqueta cuando GitHub esté disponible.
  \item Regenerar el PDF con \file{make logbook} y comprobar visualmente sus páginas.
  \item Mantener el archivo fuente \file{bitacora.tex} y el PDF generado bajo control de versiones.
\end{itemize}

\section{Resumen del alcance vigente}

El alcance vigente define una prueba de concepto académica con una fuente activa SQL Server y AdventureWorks en modo de sólo lectura. La plataforma extrae metadatos, solicita propuestas estructuradas a un LLM, exige aprobación humana y validaciones determinísticas, ejecuta un ETL controlado hacia PostgreSQL y publica KPI variables, visualizaciones, hallazgos explicables, conversación contextual y reportes PDF/Excel.

El 25 de septiembre de 2026 el tutor retiró pronóstico, MAPE y RMSE porque constituyen un problema predictivo distinto de la construcción semiautomatizada del datamart. Los módulos de seguridad, parámetros, negocio y reportes están implementados con RBAC, auditoría y separación por perfiles. El núcleo funcional se encuentra completo; permanecen el cierre académico, usabilidad formal, juicio de expertos y contraste secundario documentado.

\section{Registro cronológico}

\subsection{24 de agosto de 2026 - Inicio del repositorio y ambiente}

\begin{longtable}{@{}p{0.20\textwidth}p{0.18\textwidth}p{0.54\textwidth}@{}}
  \toprule
  \textbf{Actividad} & \textbf{Estado} & \textbf{Resultado} \\
  \midrule
  \endfirsthead
  \toprule
  \textbf{Actividad} & \textbf{Estado} & \textbf{Resultado} \\
  \midrule
  \endhead
  Revisión del anteproyecto & \statusok & Se inspeccionaron las 11 páginas, incluyendo alcance, exclusiones, objetivos, metodología, cronograma y requisitos mínimos institucionales. \\
  Arquitectura base & \statusok & Se fijaron React con Vite y TypeScript, FastAPI, PostgreSQL interno/datamart y SQL Server AdventureWorks como fuente externa. \\
  Docker-first & \statusok & El host sólo requiere Docker Compose. Se definieron etapas de desarrollo, prueba y producción para frontend y backend. \\
  Bases reproducibles & \statusok & Se añadieron imágenes inicializadoras para PostgreSQL y SQL Server. AdventureWorks se restaura desde el respaldo oficial en un volumen vacío. \\
  Seguridad & \statusok & El paquete \file{security} y las reglas del RBAC quedaron documentados, sin modelos ni endpoints funcionales. \\
  Observabilidad & \statusok & Se añadieron endpoints \file{/health/live} y \file{/health/ready}, más una pantalla de estado en React. \\
  DevOps & \statusok & Se prepararon CI, CD, Dependabot, plantilla de pull request y publicación de imágenes en GHCR. \\
  Guía de réplica & \statusok & Se documentaron el entorno aislado desde cero y el consumo de imágenes publicadas. \\
  GitHub remoto & \statusok & Repositorio privado \file{LFXAS/prototipo-bi-ia}, rama \file{main} publicada y remoto \file{origin} configurado. \\
  Verificación completa & \statusok & El conjunto se reconstruyó desde volúmenes vacíos; los cuatro servicios quedaron saludables y las pruebas automatizadas finalizaron correctamente. \\
  \bottomrule
\end{longtable}

\subsubsection{Decisiones técnicas}

\begin{itemize}
  \item El repositorio definitivo se ubica en \file{/home/ceo/Proyectos/BI}; no se codifican rutas absolutas en Compose, por lo que otras máquinas pueden clonarlo en cualquier ubicación.
  \item El Compose predeterminado crea red, bases y volúmenes aislados. Usa los puertos anfitrión 55432 y 51433 para no interferir con instancias existentes.
  \item Los volúmenes Docker no se exportan como imágenes. El estado inicial se reconstruye con respaldos públicos, scripts y futuras migraciones versionadas.
  \item \file{ApplicationIntent=ReadOnly} no es un control suficiente. El usuario \file{bi\_reader} recibe lectura y denegaciones explícitas de escritura en SQL Server.
  \item CD publicará las imágenes propias en GitHub Container Registry con etiquetas de rama, versión y SHA.
\end{itemize}

\subsubsection{Archivos y componentes principales}

\begin{itemize}
  \item \file{compose.yaml}: aplicación, PostgreSQL y SQL Server autocontenidos.
  \item \file{compose.release.yaml}: ejecución sin compilar, usando GHCR.
  \item \file{backend/}: API, configuración, sesión PostgreSQL, salud, Alembic y scaffolding modular.
  \item \file{frontend/}: interfaz React, proxy hacia la API, pruebas y Nginx de producción.
  \item \file{infra/}: inicialización de PostgreSQL, restauración de SQL Server y compilador LaTeX.
  \item \file{docs/replication-guide.md}: manual operativo paso a paso.
  \item \file{.github/workflows/}: integración y entrega continua.
\end{itemize}

\subsubsection{Incidentes y observaciones}

\begin{enumerate}
  \item La carpeta generada inicialmente contenía un directorio \file{.git} vacío y no era todavía un repositorio válido. Se resolvió inicializando Git sobre \file{main}; el commit base es \file{ead11fb}.
  \item El primer intento de generar \file{package-lock.json} falló porque npm trató de escribir en una caché de sólo lectura dentro del directorio personal. Se resolvió con la caché temporal \path{/tmp/bi-npm-cache}; el lockfile quedó generado sin vulnerabilidades reportadas.
  \item Una verificación de Compose se lanzó desde la subcarpeta \file{frontend}; no encontró los archivos raíz. Se repitió desde la raíz y las configuraciones de desarrollo y publicación resultaron válidas.
  \item La primera construcción del backend detectó formato pendiente y un import moderno requerido por Ruff. Se corrigieron ambos; Ruff, Mypy y dos pruebas Pytest finalizaron correctamente.
  \item La primera prueba del frontend duplicó la inicialización de \file{jest-dom}. Se retiró la importación redundante; ESLint, Vitest y Vite finalizaron correctamente.
  \item npm detectó vulnerabilidades en las versiones iniciales de Vite y Vitest. Se actualizaron a las versiones corregidas 7.3.6 y 3.2.7; el lockfile registra cero vulnerabilidades.
  \item La primera salud del frontend consultó \file{localhost}, que BusyBox resolvió como IPv6 mientras Vite escuchaba en IPv4. Se cambió la sonda a \file{127.0.0.1} y el contenedor quedó saludable.
  \item La primera restauración creó \file{bi\_reader}, pero una denegación demasiado amplia de \file{CONTROL} también bloqueó el permiso subordinado de conexión. Se eliminó esa denegación, la sonda de salud pasó a probar el usuario real y se reconstruyó todo desde volúmenes vacíos.
  \item La primera restauración de AdventureWorks requiere red y puede tardar varios minutos. Los reinicios posteriores reutilizan el volumen persistente.
  \item La primera ejecución de Actions finalizó correctamente, pero advirtió que varias acciones todavía declaraban Node.js 20. Se actualizaron a las versiones oficiales vigentes: \file{checkout@v7}, \file{upload-artifact@v7} y acciones Docker de generaciones 4, 6 y 7.
  \item La búsqueda autorizada de tres correos no devolvió cuentas públicas y la API de colaboradores exige nombres de usuario. Las invitaciones quedan pendientes hasta obtener los identificadores GitHub o iniciar sesión en la interfaz web.
\end{enumerate}

\subsection{24 de agosto de 2026 - Formalización del Sprint 1 y manual técnico}

\begin{itemize}
  \item Se formalizó la fase de ambiente como \textbf{Sprint 1: ambiente reproducible, repositorio y DevOps}, con objetivo, backlog, criterios de aceptación, revisión, incidencias y retrospectiva.
  \item Se creó el informe académico \file{docs/sprints/sprint-01-entorno.tex} y su PDF, acompañado por capturas reales del frontend y OpenAPI.
  \item Se creó el manual \file{docs/manual-tecnico/manual-tecnico.tex} y su PDF con instalación, configuración, réplica, operación, pruebas, CI/CD, seguridad, diagnóstico y recuperación.
  \item \file{make docs} regenera los tres documentos con la imagen LaTeX del proyecto; \file{make verify} los incluye en la definición de terminado.
  \item CI compila, compara y publica los tres PDF como un único artefacto de documentación técnica.
\end{itemize}

\subsection{24 de agosto de 2026 - Gobierno Git, colaboradores y estabilidad DevOps}

\begin{longtable}{@{}p{0.22\textwidth}p{0.18\textwidth}p{0.52\textwidth}@{}}
  \toprule
  \textbf{Actividad} & \textbf{Estado} & \textbf{Resultado} \\
  \midrule
  \endfirsthead
  \toprule
  \textbf{Actividad} & \textbf{Estado} & \textbf{Resultado} \\
  \midrule
  \endhead
  Colaboradores GitHub & \statusok & Se verificaron \file{LeslyVelasquez} y \file{EmyRoblesBerry}; ambas cuentas tienen permiso de escritura en \file{LFXAS/prototipo-bi-ia}. \\
  Ramas permanentes & \statusok & \file{develop} quedó como rama predeterminada de integración y \file{main} como rama estable de publicación. El trabajo usa \file{feature/*}, \file{fix/*}, \file{docs/*} o \file{chore/*}. \\
  Visibilidad & \statusok & Por autorización expresa, el repositorio cambió de privado a público para habilitar las reglas de protección sin contratar todavía GitHub Pro. La visibilidad de los paquetes GHCR se gestiona por separado. \\
  Protección de ramas & \statusok & \file{develop} y \file{main} exigen pull request, una aprobación, conversaciones resueltas, rama actualizada y ocho controles de CI; force-push y eliminación están deshabilitados. \\
  Flujo automático & \statusok & CI valida ambas ramas y rechaza promociones a \file{main} cuyo origen no sea \file{develop}; CD continúa limitado a \file{main} y etiquetas \file{v*}. \\
  Dependabot & \statusok & Las ramas \file{dependabot/*} se identificaron como temporales. npm agrupa sólo parches y cambios menores; los saltos mayores quedan fuera de las actualizaciones rutinarias. \\
  Reinicio SQL Server & \statusok & El inicializador espera que AdventureWorks esté en línea antes de configurar el lector y crear la marca de salud. Un reinicio con volumen existente terminó con cuatro servicios saludables y 71 tablas. \\
  Manual y guía & \statusok & README, guía de réplica, flujo Git y manual técnico describen comandos, promociones, protecciones, Dependabot y recuperación. \\
  Verificación final & \statusok & \file{make verify} validó Compose, política de ramas, backend, frontend y los tres PDF; Actionlint 1.7.12 validó los flujos, todo dentro de Docker. \\
  \bottomrule
\end{longtable}

\subsubsection{Incidentes y resolución}

\begin{enumerate}
  \item El PR automático del frontend agrupó 18 actualizaciones y falló durante \file{npm ci}. Se reprodujo el error: TypeScript 7.0.2 era incompatible con \file{typescript-}\allowbreak\file{eslint} 8.67.0, que admite versiones menores que 6.1. Se corrigió la política para no agrupar ni proponer saltos mayores rutinarios; el PR incompatible debe cerrarse sin fusionar.
  \item La captura de Actions también mostraba avisos por acciones basadas en Node.js 20. La rama \file{main} ya contenía las generaciones vigentes de \file{checkout}, \file{upload-artifact} y las acciones Docker; los avisos procedían de una rama Dependabot atrasada dos commits.
  \item El primer intento de proteger el repositorio privado devolvió HTTP 403 porque el plan actual no incluye protección de ramas privadas. El usuario autorizó volver público el repositorio; la misma configuración se aplicó correctamente a \file{develop} y \file{main}.
  \item Una prueba adicional \file{make lint} se detuvo porque SQL Server estaba no saludable después de un reinicio. El motor aceptaba conexiones, pero AdventureWorks seguía recuperándose cuando el inicializador intentó configurarla y dejó ausente la marca de disponibilidad. Se añadió una espera explícita a la base, se recreó el contenedor sin borrar el volumen y la prueba de reinicio pasó.
\end{enumerate}

\subsection{24 de agosto de 2026 - Ambientes simultáneos y vista de entrega}

\begin{longtable}{@{}p{0.22\textwidth}p{0.18\textwidth}p{0.52\textwidth}@{}}
  \toprule
  \textbf{Actividad} & \textbf{Estado} & \textbf{Resultado} \\
  \midrule
  \endfirsthead
  \toprule
  \textbf{Actividad} & \textbf{Estado} & \textbf{Resultado} \\
  \midrule
  \endhead
  Diagnóstico de acceso & \statusok & Se comprobó que Vite respondía en 5173 y FastAPI estaba disponible en 8000; el error del navegador en 8080 se debía a que todavía no existía un servicio de entrega activo en ese puerto. \\
  Vista Nginx paralela & \statusok & El perfil Compose \file{delivery} añade \file{frontend-delivery} en 8080, comparte la API y las bases de desarrollo y conserva Vite en 5173. \\
  Entrega completa aislada & \statusok & \file{compose.release.yaml} usa el proyecto \file{bi-ia-prototype-release}, API 18000, PostgreSQL 55433, SQL Server 51434 y volúmenes independientes; no reemplaza desarrollo. \\
  Automatización & \statusok & \file{make delivery-preview-up} inicia la vista ligera. \file{make release-up} exige \file{.env.release} y levanta las imágenes GHCR; \file{make release-down} detiene sólo ese ambiente. \\
  CI/CD & \statusok & CI valida desarrollo, perfil de entrega y Compose GHCR mediante contenedores efímeros. CD publica desde \file{main} o etiquetas \file{v*}; el despliegue selecciona una etiqueta sin compartir volúmenes. \\
  Guías & \statusok & README, arquitectura, réplica, flujo Git y manual técnico distinguen ramas, ambientes y consumo de recursos. \\
  Verificación local & \statusok & Cinco servicios quedaron saludables; Vite respondió en 5173, Nginx en 8080 y el proxy de entrega devolvió salud de FastAPI. \\
  Nube académica & \statuspending & Se recomienda Azure for Students con una VM Linux x64 de 2 vCPU y 8 GiB, OIDC, aprobación del ambiente y apagado automático. Falta activar la suscripción antes de crear recursos. \\
  Visibilidad GHCR & \statusok & Con autorización expresa se cambió a pública la visibilidad de \file{prototipo-bi-ia-frontend}, \file{prototipo-bi-ia-backend}, \file{prototipo-bi-ia-postgres}, \file{prototipo-bi-ia-sqlserver} y \file{prototipo-bi-ia-docs}. GitHub confirmó el estado público de cada paquete. \\
  \bottomrule
\end{longtable}

\subsubsection{Decisión operativa}

La vista Nginx es la opción recomendada para revisar rápidamente el frontend compilado porque no duplica SQL Server. La entrega GHCR completa se reserva para validar aislamiento, imágenes y restauración reproducible; puede convivir con desarrollo si el computador dispone de memoria suficiente. La vista ligera y la entrega completa comparten por defecto el puerto anfitrión 8080, por lo que se usa sólo una de ellas a la vez o se cambia su variable de puerto.

\subsubsection{Incidente de verificación documental}

La primera regeneración mostró avisos porque TeX intentó crear su caché de fuentes en una ruta sin escritura. La compilación usó su alternativa temporal y \file{make verify} terminó correctamente, pero se fijaron \file{HOME}, \file{TEXMFVAR} y \file{VARTEXFONTS} a rutas temporales dentro del contenedor para eliminar la dependencia implícita y mantener iguales los comandos locales y de CI.

La primera CI del commit \file{7f39517} falló únicamente en \file{Validate Compose}: \file{.env.release.example} existía localmente, pero una regla amplia de \file{.gitignore} impedía versionarlo. Se añadió una excepción explícita, se incorporó el archivo de ejemplo sin credenciales reales y se repitió la verificación. Este hallazgo confirma el valor de probar desde un clon limpio además del entorno local.

\clearpage
\subsection{24 de agosto de 2026 - Publicación abierta de imágenes GHCR}

\begin{longtable}{@{}p{0.22\textwidth}p{0.18\textwidth}p{0.52\textwidth}@{}}
  \toprule
  \textbf{Actividad} & \textbf{Estado} & \textbf{Resultado} \\
  \midrule
  \endfirsthead
  \toprule
  \textbf{Actividad} & \textbf{Estado} & \textbf{Resultado} \\
  \midrule
  \endhead
  Autorización & \statusok & El usuario autorizó publicar las cinco imágenes porque el repositorio es académico y las colaboradoras deben poder reproducir el entorno. No se publicaron secretos, volúmenes ni respaldos privados. \\
  Paquetes públicos & \statusok & Las cinco configuraciones de GitHub Packages indican que el paquete está público; la vista del perfil muestra los cinco contenedores sin la etiqueta \file{Private}. \\
  Acceso del equipo & \statusok & \file{LeslyVelasquez} y \file{EmyRoblesBerry} conservan permiso de escritura heredado del repositorio. La lectura de los paquetes públicos no depende de ser colaborador. \\
  Pull request 12 & \statuspending & Los ocho controles de CI están aprobados y la fusión automática está activa. Falta una aprobación de cualquiera de las dos colaboradoras para integrar los tres commits de \file{chore/git-flow} en \file{develop}. \\
  Réplica externa & \statuspending & Falta ejecutar \file{make release-up} desde un clon limpio en un segundo computador sin sesión de GHCR y conservar como evidencia las sondas de salud y los digest descargados. \\
  \bottomrule
\end{longtable}

\subsubsection{Incidentes y resolución}

\begin{enumerate}
  \item La primera prueba anónima devolvió HTTP 401 porque la visibilidad del repositorio y la de GitHub Packages son controles independientes. Se corrigió manualmente la visibilidad de cada paquete desde la cuenta propietaria.
  \item La ampliación de permisos de GitHub CLI fue autorizada mediante el flujo de dispositivo, pero el entorno protegido no pudo escribir temporalmente su archivo local de autenticación. Se completó la operación desde la sesión web autenticada, sin registrar contraseñas, códigos ni tokens en el repositorio.
  \item El mismo bloqueo de escritura impidió actualizar la bitácora durante esa intervención. No hubo pérdida ni corrupción: el árbol Git permaneció limpio y la versión 0.1.7 registró la resolución en una sesión posterior.
\end{enumerate}

\subsection{24 de agosto de 2026 - Operación documentada de contenedores}

\begin{longtable}{@{}p{0.22\textwidth}p{0.18\textwidth}p{0.52\textwidth}@{}}
  \toprule
  \textbf{Actividad} & \textbf{Estado} & \textbf{Resultado} \\
  \midrule
  \endfirsthead
  \toprule
  \textbf{Actividad} & \textbf{Estado} & \textbf{Resultado} \\
  \midrule
  \endhead
  Manual técnico & \statusok & Se elevó a la versión 1.4.0 y se añadió una sección operativa completa para los tres ambientes Docker. \\
  Desarrollo & \statusok & Se documentaron primera ejecución, construcción, pausa con \file{stop}, reanudación con \file{start} y retiro con \file{make down}. \\
  Vista Nginx & \statusok & Se documentaron inicio, consulta, pausa, reanudación y retiro exclusivo de \file{frontend-delivery}, sin afectar los cuatro servicios de desarrollo. \\
  Entrega GHCR & \statusok & Se documentaron \file{release-up}, pausa, reanudación, consulta y \file{release-down} con \file{.env.release} y Compose aislado. \\
  Persistencia & \statusok & El manual diferencia \file{stop}, \file{start}, \file{down} y \file{down --volumes}; la eliminación de volúmenes queda marcada como destructiva. \\
  Verificación & \statusok & Los PDF del manual y la bitácora se regeneraron e inspeccionaron; \file{make verify} finalizó correctamente dentro de Docker. \\
  \bottomrule
\end{longtable}

\clearpage
\subsection{1 de septiembre de 2026 - Cierre documental y adopción de SDD}

\begin{longtable}{@{}p{0.22\textwidth}p{0.18\textwidth}p{0.52\textwidth}@{}}
  \toprule
  \textbf{Actividad} & \textbf{Estado} & \textbf{Resultado} \\
  \midrule
  \endfirsthead
  \toprule
  \textbf{Actividad} & \textbf{Estado} & \textbf{Resultado} \\
  \midrule
  \endhead
  Auditoría documental & \statusok & Se revisaron README, arquitectura, réplica, flujo Git, manual técnico, informe y bitácora. Se corrigió el lenguaje heredado que todavía describía repositorio o imágenes como privados; el estado vigente es repositorio público, paquetes GHCR públicos y escritura limitada a colaboradoras autorizadas. \\
  Manual técnico & \statusok & Se actualizó a versión 1.5.0 con explicación paso a paso de CI, CD, DevOps, lectura de GitHub Actions, límites actuales del despliegue y responsabilidades del equipo. \\
  SDD & \statusok & Se incorporó la guía \file{docs/sdd-workflow.md}, índice, plantilla y la especificación inicial \file{SPR-02-rbac-y-parametros.md}. A partir del Sprint 2 toda capacidad funcional debe tener una especificación versionada antes del PR de implementación. \\
  Sprint 1 & \statusok & El informe se actualizó para reconocer la evidencia de réplica y el cierre documental sin declarar implementado RBAC, ETL, IA ni despliegue Azure. \\
  Entrega Git & \statuspending & Los cambios documentales deben seguir el flujo normal: rama \file{docs/*}, PR hacia \file{develop}, CI verde, una aprobación y posterior promoción \file{develop -> main}. \\
  \bottomrule
\end{longtable}

\subsubsection{Decisión de gobierno}

SDD se adopta como norma de trabajo, no como una promesa de implementación automática. La especificación define alcance y criterios de aceptación; Docker y CI aportan la evidencia técnica; la bitácora registra el resultado; y el pull request conserva la aprobación. Las modificaciones meramente editoriales pueden justificar su alcance en el PR, pero todo cambio funcional, de datos, seguridad, IA o integración externa requiere especificación previa.

\subsection{2 de septiembre de 2026 - Profundización documental y material docente anónimo}

\begin{longtable}{@{}p{0.22\textwidth}p{0.18\textwidth}p{0.52\textwidth}@{}}
  \toprule
  \textbf{Actividad} & \textbf{Estado} & \textbf{Resultado} \\
  \midrule
  \endfirsthead
  \toprule
  \textbf{Actividad} & \textbf{Estado} & \textbf{Resultado} \\
  \midrule
  \endhead
  Estado GitHub & \statusok & Se verificó que el PR \#15 fue fusionado en \file{develop} y el PR \#16 promovió el cierre del Sprint 1 a \file{main}. \\
  Manual técnico & \statusok & Se elevó a versión 1.7.0 y se incorporó una reconstrucción operativa desde cero, con explicación de variables, Compose, Dockerfiles, Makefile, scripts, Git, GitHub, CI, CD, GHCR y persistencia, además de capturas reales de Actions y los controles de promoción. \\
  Informe Sprint 1 & \statusok & Se corrigió el estado de cierre y se agregó una matriz de trazabilidad entre necesidades, archivos, controles y evidencias, además del ciclo DevOps aplicado. \\
  Manual docente & \statusok & Se preparó fuera del repositorio público un documento LaTeX anónimo con un caso genérico, ruta manual completa, ruta asistida por IA y diagramas de arquitectura y CI. No contiene nombres, cuentas, rutas ni identificadores del proyecto real. \\
  Verificación & \statusok & Se compilaron los cuatro PDF, se renderizaron e inspeccionaron todas sus páginas y \file{make verify} validó Compose, ramas, backend, frontend y documentación dentro de Docker. \\
  Entrega Git & \statusok & La mejora se consolidó y publicó en la rama \file{docs/ampliar-manuales-sprint-1}; el PR \#17 se abrió hacia \file{develop} bajo el flujo protegido después de la verificación local. \\
  \bottomrule
\end{longtable}

\subsubsection{Criterio de separación documental}

El manual técnico y el informe del sprint pertenecen a la tesis y conservan referencias verificables al repositorio real. El manual docente es un artefacto independiente y anónimo: enseña el mismo tipo de proceso mediante nombres ficticios y marcadores, sin publicar información que permita asociarlo con las estudiantes o con el prototipo.

\clearpage
\subsection{4 de septiembre de 2026 - Especificación UX y configuración LLM del Sprint 2}

\begin{longtable}{@{}p{0.22\textwidth}p{0.18\textwidth}p{0.52\textwidth}@{}}
  \toprule
  \textbf{Actividad} & \textbf{Estado} & \textbf{Resultado} \\
  \midrule
  \endfirsthead
  \toprule
  \textbf{Actividad} & \textbf{Estado} & \textbf{Resultado} \\
  \midrule
  \endhead
  Revisión SDD & \statusok & La especificación \file{SPR-02-rbac-y-parametros.md} se amplió antes de escribir código. El alcance conserva RBAC y parámetros, e incorpora la experiencia UI/UX como requisito funcional verificable. \\
  Adaptación responsive & \statusok & Se definieron contextos de prueba para móvil desde 320 px, tableta, escritorio y pantalla amplia. Se prohíbe depender de anchos fijos o generar desplazamiento horizontal involuntario. \\
  Navegación y flujos & \statusok & Se documentaron los flujos de acceso, navegación autorizada, administración de usuarios, roles/permisos y parámetros/conexión, incluidos los estados de sesión vencida y acceso denegado. \\
  Componentes y accesibilidad & \statusok & Se acordaron cascarón de aplicación, formularios, tablas, diálogos, avisos y pantallas de estado, con foco visible, etiquetas, mensajes claros y contraste verificable. \\
  Plantilla SDD & \statusok & La plantilla y el índice de especificaciones ahora exigen declarar UI/UX responsive y accesibilidad para toda capacidad futura que exponga interfaz. \\
  Configuración LLM & \statusok & Se aprobó una única configuración activa entre Gemini Cloud, Qwen Cloud u Ollama local. Sus referencias de secreto son \file{GEMINI\_API\_KEY}, \file{DASHSCOPE\_API\_KEY} y \file{none}, respectivamente; se conservan fuera de PostgreSQL y de la interfaz. Una prueba real y acotada confirmará conectividad sin enviar datos del negocio. \\
  ADR de proveedores & \statusok & El ADR 0002 fija adaptadores internos de FastAPI, un catálogo cerrado de tres proveedores y \file{qwen3:4b} como modelo local recomendado. Añadir otro proveedor requerirá actualizar decisión, especificación, controles de secretos y pruebas. \\
  Implementación & \statuspending & No se implementaron módulos ni pantallas en esta actividad. La especificación debe ser revisada y aprobada antes de crear la rama de implementación \file{feature/rbac-y-parametros}. \\
  \bottomrule
\end{longtable}

\subsubsection{Decisión de diseño}

La plataforma usará un cascarón React reusable y menús derivados de permisos ya validados por FastAPI. La interfaz ofrece orientación y prevención de errores, pero no reemplaza las decisiones de autorización del backend. Se privilegiarán componentes nativos y CSS mantenible; una biblioteca visual externa sólo se evaluará con justificación explícita. La configuración del LLM será administrable con un único proveedor activo: Gemini Cloud, Qwen Cloud u Ollama local. Su conexión será comprobable mediante una prueba real, limitada y sin datos del negocio. Las solicitudes generativas sobre BI sólo podrán ser realizadas posteriormente por el módulo \file{copilot}, usando una configuración activa validada y sin exponer credenciales.

\clearpage
\subsection{6 de septiembre de 2026 - Implementación inicial del Sprint 2}

\begin{longtable}{@{}p{0.22\textwidth}p{0.18\textwidth}p{0.52\textwidth}@{}}
  \toprule
  \textbf{Actividad} & \textbf{Estado} & \textbf{Resultado} \\
  \midrule
  \endfirsthead
  \toprule
  \textbf{Actividad} & \textbf{Estado} & \textbf{Resultado} \\
  \midrule
  \endhead
  Migración Sprint 2 & \statusok & Se incorporaron las revisiones Alembic \file{20260906\_01}, \file{20260906\_02} y \file{20260906\_03}, que crean las tablas iniciales y añaden ciclo de vida de permisos y agrupación administrable de menús. El servicio \file{migrate} las ejecuta antes de FastAPI. \\
  Seguridad & \statusok & Se implementaron hash Argon2, JWT de vida limitada, cuenta administradora inicial por variables de entorno, permisos estables, menús autorizados y auditoría de acceso y cambios administrativos. \\
  Parámetros y fuente externa & \statusok & Se implementaron parámetros no secretos, prueba lectora de AdventureWorks mediante la cuenta configurada y adaptadores acotados para Gemini, Qwen Cloud y Ollama local. Las claves siguen exclusivamente en variables de entorno. \\
  Interfaz administrativa & \statusok & React incorpora acceso, cierre de sesión, navegación derivada de permisos y agrupada por módulo, menú hamburguesa en tableta y móvil, avisos, formularios de alta y edición, asociaciones por identificadores y activación/desactivación lógica. Las listas de usuarios, roles, permisos, menús, parámetros, LLM y auditoría reciben paginación desde API. La auditoría queda exclusivamente en consulta. \\
  Verificación integrada & \statusok & En Docker se comprobó la migración, la creación del administrador, el inicio de sesión JWT, la sesión con ocho menús autorizados y las consultas protegidas de usuarios, roles, permisos, menús, parámetros, LLM y auditoría. \\
  Pendiente de cierre & \statuspending & Faltan pruebas negativas de autorización y de adaptadores LLM, evidencia visual responsive documentada, \file{make verify} final y el PR hacia \file{develop}. \\
  Alcance CRUD & \statusok & La especificación se precisó para exigir CRUD completo de usuarios, roles, permisos, menús, parámetros y configuración LLM, con desactivación segura, autorización backend y auditoría. \\
  \bottomrule
\end{longtable}

\clearpage
\subsection{10 de septiembre de 2026 - Corrección funcional y visual del módulo administrativo}

\begin{longtable}{@{}p{0.22\textwidth}p{0.18\textwidth}p{0.52\textwidth}@{}}
  \toprule
  \textbf{Actividad} & \textbf{Estado} & \textbf{Resultado} \\
  \midrule
  \endfirsthead
  \toprule
  \textbf{Actividad} & \textbf{Estado} & \textbf{Resultado} \\
  \midrule
  \endhead
  Especificación corregida & \statusok & Las observaciones funcionales se reorganizaron en especificaciones de seguridad RBAC, navegación, parámetros/LLM, identidad visual y una matriz de trazabilidad. No se adelantan ETL, analítica ni Copiloto BI. \\
  Protección administrativa & \statusok & La migración \file{20260910\_04} añade la marca persistente \file{is\_system\_protected}. La cuenta y el rol administrativo inicial no se pueden desactivar desde la API; el rol debe conservar permisos de recuperación y los permisos/menús de sistema no se desactivan. \\
  Interfaz RBAC & \statusok & La creación y edición de usuarios selecciona roles por nombre; los roles seleccionan permisos por nombre y descripción. La interfaz ya no solicita identificadores numéricos ni códigos técnicos para esas asociaciones. \\
  Menús y estado de pantalla & \statusok & Inicio se muestra como acceso directo, sin grupo \emph{Principal}. Los grupos pueden plegarse aun cuando contienen la ruta activa, el sidebar se puede ocultar en escritorio y cada ruta reinicia su estado de edición para evitar datos cruzados o valores \file{undefined}. \\
  Ajuste de controles administrativos & \statusok & Los listbox de permisos otorgados y requeridos ocupan la fila completa y se adaptan a la opción legible más larga dentro de límites responsive. El control textual del encabezado se sustituyó por un chevrón compacto superpuesto al borde del sidebar, visible también cuando la barra está oculta y sin desplazar la navegación. \\
  Mensajes de validación & \statusok & La interfaz traduce el detalle estructurado de validación de FastAPI para evitar el texto \file{[object Object]}. El correo de creación de usuario se solicita mediante un control de correo y presenta un ejemplo de formato válido. \\
  Respuesta de usuario creado & \statusok & Se corrigió la carga diferida de permisos al devolver un usuario recién creado o editado. La API vuelve a consultar el usuario con sus roles y permisos antes de responder, evitando un error 500 después de confirmar la transacción. \\
  Parámetros y LLM & \statusok & Parámetros generales no permite claves libres mientras no exista un catálogo consumido por un módulo. Configuración LLM presenta proveedor guiado, referencia de secreto por entorno y acción para probar la conexión desde FastAPI. \\
  Dirección visual & \statusok & Se definieron tokens CSS para marca, superficies y estados, además del cascarón futuro para Datos, Analítica e IA. Los módulos futuros no se muestran como opciones vacías en Sprint 2. \\
  Verificación & \statusok & Migración aplicada en Docker; backend saludable; Ruff, Mypy y cinco pruebas Pytest aprobadas; ESLint, dos pruebas Vitest y construcción Vite aprobados. \file{make verify} aprobó configuración Compose, flujo de ramas, imágenes de prueba y generación de PDF. \\
  \bottomrule
\end{longtable}

\clearpage
\subsection{11 de septiembre de 2026 - Alternativa local Ollama para configuración LLM}

\begin{longtable}{@{}p{0.22\textwidth}p{0.18\textwidth}p{0.52\textwidth}@{}}
  \toprule
  \textbf{Actividad} & \textbf{Estado} & \textbf{Resultado} \\
  \midrule
  \endfirsthead
  \toprule
  \textbf{Actividad} & \textbf{Estado} & \textbf{Resultado} \\
  \midrule
  \endhead
  Especificación y decisión & \statusok & La especificación SPR-02-03, la matriz de observaciones y el ADR 0002 precisan que Ollama es una alternativa local opcional del Sprint 2. No modifica el alcance histórico del Sprint 1 ni adelanta el módulo Copiloto BI. \\
  Perfil Docker & \statusok & \file{compose.yaml} incorpora el perfil \file{local-llm} con el servicio \file{ollama}, volumen nombrado \file{ollama\_models}, sin publicar el puerto 11434 en el host. \\
  Operación por equipo & \statusok & \file{make ollama-up}, \file{make ollama-pull}, \file{make ollama-status} y \file{make ollama-down} permiten iniciar, descargar, revisar y detener el servicio. El modelo se conserva por computador y no se versiona, publica en GHCR ni descarga en CI/CD. \\
  Validación segura & \statusok & FastAPI consulta \file{/api/tags} y exige que el modelo configurado esté descargado antes de informar éxito. Una prueba unitaria cubre modelo presente y ausente; la prueba real confirmó \file{qwen3:4b} sin enviar datos del negocio. \\
  Configuración funcional & \statusok & Se registró una configuración local de prueba inactiva, de modo que no se sustituyó la configuración Gemini existente. La creación y prueba conservaron la auditoría sin persistir claves. \\
  Documentación & \statusok & README, guía de réplica, arquitectura y manual técnico se ampliaron con requisitos, comandos, persistencia, URL interna y procedimiento para cada programadora. Se creó el informe formal del Sprint 2 y CI compila/verifica ahora los cuatro PDF oficiales. \\
  \bottomrule
\end{longtable}

\clearpage
\subsection{11 de septiembre de 2026 - Corrección de auditoría y cuentas temporales}

\begin{longtable}{@{}p{0.22\textwidth}p{0.18\textwidth}p{0.52\textwidth}@{}}
  \toprule
  \textbf{Actividad} & \textbf{Estado} & \textbf{Resultado} \\
  \midrule
  \endfirsthead
  \toprule
  \textbf{Actividad} & \textbf{Estado} & \textbf{Resultado} \\
  \midrule
  \endhead
  Diagnóstico & \statusok & Se comprobó que una clave foránea restrictiva desde \file{audit\_events.actor\_user\_id} impedía eliminar usuarios temporales que habían ejecutado acciones. La regla era demasiado restrictiva para cuentas de prueba. \\
  Migración & \statusok & La revisión Alembic \file{20260911\_05} incorpora \file{actor\_label}, conserva la identidad histórica segura y cambia la relación a \file{ON DELETE SET NULL}. \\
  Regla funcional & \statusok & Una cuenta no protegida puede eliminarse después de resolver sus asociaciones usuario--rol. Los eventos de auditoría no se eliminan: conservan etiqueta histórica y su referencia técnica al actor queda vacía. \\
  Evidencia & \statusok & Una prueba transaccional creó y eliminó una cuenta temporal; el evento sobrevivió con \file{actor\_user\_id} nulo y etiqueta histórica. La transacción se revirtió sin dejar datos de prueba. \\
  Trazabilidad & \statusok & La especificación RBAC y la matriz SDD incorporan OBS-21. El cambio queda en el commit \file{ec3cdea} y PR \#30 hacia \file{develop}. \\
  \bottomrule
\end{longtable}

\clearpage
\subsection{11 de septiembre de 2026 - Migraciones y catálogo base reproducible}

\begin{longtable}{@{}p{0.22\textwidth}p{0.18\textwidth}p{0.52\textwidth}@{}}
  \toprule
  \textbf{Actividad} & \textbf{Estado} & \textbf{Resultado} \\
  \midrule
  \endfirsthead
  \toprule
  \textbf{Actividad} & \textbf{Estado} & \textbf{Resultado} \\
  \midrule
  \endhead
  Decisión de persistencia & \statusok & Se confirma que los volúmenes PostgreSQL son estado local por equipo y no se comparten mediante GitHub ni como mecanismo de actualización de sprints. \\
  Servicio de semillas & \statusok & Compose ejecuta \file{seed} después de \file{migrate} y antes de FastAPI. El proceso es idempotente y carga el catálogo protegido de permisos, menús, rol administrador y cuenta inicial. \\
  Conservación de datos & \statusok & La semilla no borra ni distribuye cuentas temporales, auditoría, configuraciones LLM, claves, tokens ni datos de negocio. Las demostraciones futuras requerirán semillas sintéticas y opcionales aprobadas por SDD. \\
  Operación y evidencia & \statusok & Se incorpora \file{make seed}; README, arquitectura, guía de réplica, especificación RBAC, manual técnico e informe del Sprint 2 documentan el flujo de actualización por equipo. \\
  \bottomrule
\end{longtable}

\subsection{11 de septiembre de 2026 - Modelo local ágil para Ollama}

\begin{longtable}{@{}p{0.22\textwidth}p{0.18\textwidth}p{0.52\textwidth}@{}}
  \toprule
  \textbf{Actividad} & \textbf{Estado} & \textbf{Resultado} \\
  \midrule
  \endfirsthead
  \toprule
  \textbf{Actividad} & \textbf{Estado} & \textbf{Resultado} \\
  \midrule
  \endhead
  Decisión operativa & \statusok & \file{qwen2.5:3b} reemplaza a \file{qwen3:4b} como modelo local predeterminado para desarrollo y demostraciones ágiles en español. \\
  Recursos locales & \statusok & El contexto predeterminado se ajusta a 2048; el modelo anterior queda como alternativa de mayor capacidad, no como requisito de cada equipo. \\
  Reproducibilidad & \statusok & \file{.env.example}, Compose, Makefile, interfaz, especificación, ADR y guías indican el mismo valor. No se versionan modelos, volúmenes ni archivos \file{.env} personales. \\
  Verificación & \statusok & Pruebas de proveedor, configuración Compose y comprobación integral ejecutadas antes del PR. \\
  \bottomrule
\end{longtable}

\clearpage
\subsection{14 de septiembre de 2026 - Especificación SDD del Sprint 3}

\begin{longtable}{@{}p{0.22\textwidth}p{0.18\textwidth}p{0.52\textwidth}@{}}
  \toprule
  \textbf{Actividad} & \textbf{Estado} & \textbf{Resultado} \\
  \midrule
  \endfirsthead
  \toprule
  \textbf{Actividad} & \textbf{Estado} & \textbf{Resultado} \\
  \midrule
  \endhead
  Alcance de investigación & \statusok & Se registra el título vigente, los cinco objetivos y la regla de que AdventureWorks es la única fuente funcional; los datos sintéticos se reservan para pruebas controladas. \\
  Especificaciones & \statusok & Se refinan el índice SPR-03 y cuatro contratos para configuración web, conexión e introspección, propuesta BI asistida y experiencia de exploración/revisión. Todos permanecen en borrador hasta aprobación humana. \\
  Usuario objetivo & \statusok & El gerente o analista expresa una necesidad en español sin seleccionar tablas ni escribir SQL; el explorador técnico queda como consulta avanzada de sólo lectura. \\
  Separación de responsabilidades & \statusok & La aplicación conecta, normaliza, preselecciona, valida y audita; el LLM propone un documento estructurado; una persona aprueba el significado de negocio. Ninguna propuesta ejecuta SQL o ETL. \\
  Parametrización segura & \statusok & Conexiones, claves LLM y parámetros operativos tipados se administrarán desde la web; los secretos se cifran, SQL Server es el primer adaptador y el contrato admite motores futuros. \\
  Interpretación dinámica & \statusok & El LLM interpreta metadatos técnicos por bloques, propone conceptos comprensibles en español y FastAPI descarta toda tabla, columna o relación inexistente. \\
  Decisión arquitectónica & \statusok & El ADR 0003 define instantáneas inmutables, mapa semántico dinámico, contrato JSON, validadores y puerta de aprobación. Sprint 4 deberá incorporar un constructor determinístico y ejecutor backend, no consultas manuales por análisis. \\
  Límite del sprint & \statusok & Sprint 3 termina con una propuesta aprobada como entrada de Sprint 4; datamart físico, ETL, KPIs calculados, dashboard y pronóstico permanecen pendientes. \\
  Simplificación de alcance & \statusok & Se separa lo obligatorio, lo preparado y lo futuro. Sprint 3 implementará un solo motor, una sola fuente, cuatro parámetros consumidos, un recorrido de negocio de cuatro pasos y un explorador técnico de sólo lectura. \\
  Valor para usuario final & \statusok & El gerente formula una necesidad, comprende los conceptos detectados y aprueba o rechaza una propuesta sin seleccionar tablas, editar archivos, escribir SQL ni interpretar JSON. \\
  \bottomrule
\end{longtable}

\clearpage
\subsection{18 de septiembre de 2026 - Especificación UI/UX profesional del Sprint 3}

\begin{longtable}{@{}p{0.22\textwidth}p{0.18\textwidth}p{0.52\textwidth}@{}}
  \toprule
  \textbf{Actividad} & \textbf{Estado} & \textbf{Resultado} \\
  \midrule
  \endfirsthead
  \toprule
  \textbf{Actividad} & \textbf{Estado} & \textbf{Resultado} \\
  \midrule
  \endhead
  Especificación UI/UX & \statusok & Se incorpora \file{SPR-03-05} como contrato normativo para la navegación, lenguaje visual, componentes, estados, contenido, responsive y accesibilidad del Sprint 3. Permanece en borrador hasta aprobación humana. \\
  Perfiles de uso & \statusok & Se separan los recorridos del administrador, gerente o solicitante y analista BI. El asistente es el acceso principal del usuario de negocio y la trazabilidad técnica queda disponible sin convertirse en una tarea obligatoria. \\
  Usuario principal & \statusok & Se precisa que el analista BI o responsable de datos opera el asistente, resuelve ambigüedades y aprueba o rechaza. El gerente aporta la necesidad y consume posteriormente los resultados; no valida relaciones, granularidad ni el plan ETL. \\
  Flujo principal & \statusok & El Asistente de análisis utiliza cuatro pasos: necesidad de negocio, conceptos encontrados, propuesta validada y revisión humana. La experiencia diferencia contenido propuesto por IA, comprobación de la aplicación y decisión de una persona. \\
  Preparación inicial & \statusok & Se añade un wizard de cinco pasos que orquesta la configuración LLM, la fuente SQL Server, las pruebas, la activación y la instantánea. Reutiliza los módulos reales, permite continuar después y deja de ser obligatorio al completar los prerrequisitos. \\
  Drag and drop & \statusok & Se descarta un editor funcional en Sprint 3 porque trasladaría al usuario el diseño técnico y requeriría grafos, reglas, deshacer, persistencia y accesibilidad adicional. El plan ETL se explica mediante una vista visual no editable; un editor avanzado requerirá otra especificación en Sprint 4. \\
  Pantallas & \statusok & Se definen Inicio orientador, Conexiones de datos, Configuración LLM, Asistente de análisis y Explorador de esquema, con acciones, estados, mensajes y comportamiento móvil y de escritorio. \\
  Consistencia visual & \statusok & Se reutilizan los tokens y el cascarón del Sprint 2; no se incorpora otra identidad ni una biblioteca visual completa. Se especifican componentes compartidos, jerarquía, microcopy y una sola acción primaria por contexto. \\
  Accesibilidad & \statusok & Se establece WCAG 2.1 AA como objetivo para los recorridos implementados, con teclado, foco visible, etiquetas, contraste, regiones vivas, movimiento reducido y verificación a 320, 768, 1024 y 1440 píxeles. \\
  Viabilidad & \statusok & Se excluyen chat abierto, diseñador gráfico, arrastrar y soltar, grafo completo, temas por usuario y módulos futuros. Las esperas no muestran porcentajes ficticios y ninguna pantalla presenta datos simulados como resultados reales. \\
  \bottomrule
\end{longtable}

\clearpage
\subsection{18 de septiembre de 2026 - Cierre de credenciales LLM web del Sprint 2}

\begin{longtable}{@{}p{0.22\textwidth}p{0.18\textwidth}p{0.52\textwidth}@{}}
  \toprule
  \textbf{Actividad} & \textbf{Estado} & \textbf{Resultado} \\
  \midrule
  \endfirsthead
  \toprule
  \textbf{Actividad} & \textbf{Estado} & \textbf{Resultado} \\
  \midrule
  \endhead
  Migración & \statusok & La revisión Alembic \file{20260918\_06} incorpora \file{app.secrets} y una relación uno a uno opcional desde \file{llm\_configurations}; no migra valores desde archivos locales. \\
  Protección & \statusok & FastAPI cifra cada API key con cifrado autenticado. La raíz se genera con permisos restrictivos en el volumen \file{secret\_key\_data}, separado de PostgreSQL y no versionado. \\
  API y autorización & \statusok & \file{PUT /api/v1/llm-configurations/\{id\}/secret} exige \file{parameters.llm.write}, registra o reemplaza la clave y responde sólo con su estado. La auditoría no contiene credenciales. \\
  Interfaz & \statusok & Gemini y Qwen ofrecen \textbf{Registrar/Reemplazar credencial}; la tabla muestra \textbf{Credencial pendiente/configurada}. Ollama indica que no requiere API key. El formulario de contraseña se limpia al guardar o cancelar. \\
  Despliegue & \statusok & Compose dejó de inyectar \file{GEMINI\_API\_KEY} y \file{DASHSCOPE\_API\_KEY}. Desarrollo y entrega declaran el volumen persistente de clave maestra. \\
  Pruebas automáticas & \statusok & Ruff, formato, Mypy, 13 pruebas Pytest, 4 pruebas Vitest, ESLint y construcción Vite finalizaron correctamente dentro de Docker. \\
  Prueba integrada & \statusok & Una credencial temporal produjo cero coincidencias en claro en PostgreSQL, permaneció utilizable después de reiniciar el backend y su eliminación dejó cero secretos huérfanos. \\
  Evidencia visual & \statusok & Se revisó la pantalla real en ancho móvil: formulario, estado, acciones y microcopy permanecen legibles, accesibles y sin mostrar la clave. \\
  Documentación & \statusok & README, guía de réplica, arquitectura, ADR 0002, especificaciones, informe del Sprint 2 y manual técnico reflejan el nuevo flujo. \\
  \bottomrule
\end{longtable}

\subsection{18 de septiembre de 2026 - Primer incremento funcional del Sprint 3}

\begin{longtable}{@{}p{0.22\textwidth}p{0.18\textwidth}p{0.52\textwidth}@{}}
  \toprule
  \textbf{Actividad} & \textbf{Estado} & \textbf{Resultado} \\
  \midrule
  \endfirsthead
  \toprule
  \textbf{Actividad} & \textbf{Estado} & \textbf{Resultado} \\
  \midrule
  \endhead
  Migración & \statusok & La revisión Alembic \file{20260918\_07} amplía \file{app.parameters}, crea \file{app.data\_connections}, relaciona cada conexión con un secreto y refuerza una sola fuente activa mediante índice parcial. \\
  Conexión web & \statusok & React y FastAPI permiten registrar, editar, probar, activar, desactivar y eliminar una fuente SQL Server. Los campos se validan por separado y no se acepta una cadena de conexión libre. \\
  Sólo lectura & \statusok & La prueba consulta el catálogo y rechaza cuentas con membresía o privilegios de escritura. AdventureWorks con \file{bi\_reader} fue validada y activada desde la API. \\
  Parámetros & \statusok & Se habilitaron cuatro claves tipadas con nombre, explicación, módulo, rango y valor predeterminado. La pantalla permite guardar dentro del rango y restaurar. \\
  Seguridad & \statusok & La contraseña se cifra con la misma raíz local de secretos, no se carga al listar conexiones y se delimita como un solo valor ODBC aunque contenga separadores. La inspección integrada produjo cero coincidencias de credenciales en claro en PostgreSQL. \\
  Persistencia & \statusok & Después de reiniciar FastAPI permanecieron una conexión activa y el estado de prueba correcto, sin exponer el secreto al cliente. \\
  Autorización y auditoría & \statusok & Se añadieron \file{connections.read}, \file{connections.write} y \file{connections.test}; las mutaciones conservan eventos seguros sin host completo, contraseña ni cadena ODBC. \\
  Pruebas & \statusok & 18 pruebas Pytest y 6 pruebas Vitest, Ruff, formato, Mypy, ESLint y construcción Vite finalizaron correctamente dentro de Docker. \\
  Límite & \statuspending & El bloqueo por dependencias se verificará cuando la siguiente migración incorpore \file{metadata\_snapshots}; introspección, propuesta con IA y ETL no pertenecen a este incremento. \\
  \bottomrule
\end{longtable}

\subsection{18 de septiembre de 2026 - Segundo incremento funcional del Sprint 3}

\begin{longtable}{@{}p{0.22\textwidth}p{0.18\textwidth}p{0.52\textwidth}@{}}
  \toprule
  \textbf{Actividad} & \textbf{Estado} & \textbf{Resultado} \\
  \midrule
  \endfirsthead
  \toprule
  \textbf{Actividad} & \textbf{Estado} & \textbf{Resultado} \\
  \midrule
  \endhead
  Migración & \statusok & La revisión Alembic \file{20260918\_08} crea \file{app.metadata\_snapshots} con JSONB, hash, totales, actor histórico, unicidad y relaciones restrictivas para conservar trazabilidad. \\
  Introspección & \statusok & FastAPI consulta catálogos SQL Server de tablas, columnas, PK y FK con la fuente activa de sólo lectura; no extrae filas, valores, procedimientos ni código SQL. \\
  Canonicalización & \statusok & El documento se ordena por esquema, tabla, posición y relación. SHA-256 identifica el contenido y una segunda captura idéntica reutiliza la versión vigente. \\
  Evidencia real & \statusok & AdventureWorks produjo 6 esquemas, 71 tablas, 444 columnas y 90 relaciones. La repetición conservó el mismo hash de 64 caracteres y no aumentó el total de instantáneas. \\
  API & \statusok & Se incorporaron fuente activa, captura, listado, detalle, búsqueda paginada y detalle de tabla; un bloqueo asesor evita actualizaciones simultáneas por conexión. \\
  Interfaz & \statusok & Conexiones muestra la fuente activa, sus totales y \textbf{Actualizar metadatos}. \textbf{Datos > Explorador de esquema} permite buscar y revisar columnas y relaciones con diseño apilable. \\
  Seguridad & \statusok & La contraseña se descifra únicamente en memoria; respuestas, auditoría e instantánea omiten host, credencial, cadena ODBC y filas del negocio. Una conexión con capturas no puede eliminarse. \\
  Autorización & \statusok & Se añadieron \file{metadata.read} y \file{metadata.refresh}, menú protegido e incorporación idempotente al rol administrador. \\
  Pruebas & \statusok & 21 pruebas Pytest y 7 pruebas Vitest, Ruff, formato, Mypy, ESLint y construcción Vite finalizaron correctamente dentro de Docker. \\
  Límite & \statusok & El incremento no invoca al LLM, no interpreta conceptos de negocio y no crea ETL ni datamart; esas capacidades conservan sus especificaciones posteriores. \\
  \bottomrule
\end{longtable}

\subsection{19 de septiembre de 2026 - Cierre funcional del Sprint 3}

\begin{longtable}{@{}p{0.22\textwidth}p{0.18\textwidth}p{0.52\textwidth}@{}}
  \toprule
  \textbf{Actividad} & \textbf{Estado} & \textbf{Resultado} \\
  \midrule
  \endfirsthead
  \toprule
  \textbf{Actividad} & \textbf{Estado} & \textbf{Resultado} \\
  \midrule
  \endhead
  Persistencia & \statusok & Las revisiones \file{20260919\_09} a \file{20260920\_13} crean propuestas BI inmutables, registran el perfil \file{ventas}, enlazan revisiones, parametrizan el razonamiento LLM y admiten estados auditables de aprobación retirada y descarte, ligando cada intento con instantánea, entrada, proveedor, versión, resultado y decisión. \\
  Asistente & \statusok & \textbf{Asistente de datamart} incorpora un catálogo de dominios y un recorrido de cinco pasos: necesidad comercial, conceptos españoles, propuesta validada, personalización opcional y decisión supervisada. Sólo ventas está habilitado; las capacidades visibles se derivan de los metadatos. \\
  Copiloto IA & \statusok & Los adaptadores Ollama, Gemini y Qwen Cloud pueden solicitar JSON acotado. La IA interpreta metadatos y decide granularidad, hecho, medidas, dimensiones y KPIs; FastAPI completa claves, relaciones, reglas y plan ETL de forma determinística sin recibir filas ni secretos. \\
  Validación & \statusok & FastAPI comprueba referencias, alcance, columnas, agregaciones, FK, KPI, operaciones y ausencia de instrucciones ejecutables. Los errores no pueden aprobarse y las advertencias requieren confirmación. \\
  Evidencia técnica & \statusok & La plataforma delimita la \textbf{Validación estructural de la propuesta - Sprint 3}: integridad de metadatos, referencias, contrato y reproducción determinística. Las propuestas reales de Ollama y Gemini coincidieron en huella; permanecen explícitamente pendientes la conciliación OLTP--datamart, contraste DW, juicio de expertos y MAPE/RMSE. \\
  Personalización & \statusok & El analista ajusta resumen, granularidad, dimensiones, medidas, agregaciones y KPIs ya verificados. Cada ajuste exige justificación, crea una versión enlazada y vuelve a validarse sin invocar al LLM ni aceptar SQL o referencias libres. \\
  Historial & \statusok & Las versiones se filtran en el servidor por dominio y estado, aparecen inicialmente en \textbf{Lista para revisar}, disponen de paginación y pueden abrirse sin consumir nuevamente el proveedor. \\
  Seguridad & \statusok & Se añadieron permisos separados de consulta, generación y revisión, eventos de auditoría seguros y estados persistidos para fallo del proveedor, fallo de validación, aprobación y rechazo. \\
  Extensibilidad & \statusok & El lector se resuelve mediante registro de conectores y el comportamiento analítico mediante perfiles de dominio. Sólo SQL Server y ventas quedan habilitados; MySQL e inventario requieren implementación y pruebas propias. \\
  Proveedores & \statusok & Ollama validó \file{qwen2.5:3b} y produjo una propuesta real posteriormente aprobada. Tras sustituir desde la web el modelo Gemini discontinuado para usuarios nuevos, \file{gemini-3.6-flash} con razonamiento mínimo produjo una segunda propuesta real válida y lista para revisión. Ambas versiones permanecen seleccionables. \\
  Pruebas & \statusok & Ruff, formato, Mypy y 38 pruebas Pytest; ESLint, 8 pruebas Vitest y construcción Vite finalizaron correctamente. \\
  Límite & \statusok & No se ejecuta SQL, no se materializa el datamart y no se calculan KPIs. Aprobar prepara el contrato de entrada del Sprint 4. \\
  Aceptación & \statuspending & Las autoras deben recorrer el Sprint 3 completo antes de crear los pull requests de integración a \file{develop}. \\
  \bottomrule
\end{longtable}

\subsection{20 de septiembre de 2026 - Corrección semántica y depuración de propuestas}

\begin{longtable}{@{}p{0.22\textwidth}p{0.18\textwidth}p{0.52\textwidth}@{}}
  \toprule
  \textbf{Actividad} & \textbf{Estado} & \textbf{Resultado} \\
  \midrule
  Contrato semántico & \statusok & Medidas y KPI declaran función tipada; el backend comprueba columna, agregación, operación y compatibilidad antes de permitir aprobación. \\
  Mensajes & \statusok & Observaciones del proveedor, ajustes automáticos, advertencias determinísticas y decisiones excluidas se muestran por separado y con acciones concretas. \\
  Supervisión & \statusok & El analista puede excluir o reasignar un KPI a una medida compatible verificada sin SQL; la ausencia de alternativas exige retirar el KPI o generar otra versión. \\
  Depuración & \statusok & Una aprobación puede retirarse y las versiones no aprobadas pueden descartarse con motivo; se conserva linaje y auditoría sin borrado físico. \\
  Compatibilidad futura & \statusok & La validación usa funciones semánticas y reglas del perfil, no excepciones ligadas a una pantalla o tabla específica. Sprint 4 deberá revalidar inmediatamente antes del ETL. \\
  Pruebas & \statusok & Ruff, Mypy, 42 pruebas Pytest, ESLint, 8 pruebas Vitest y construcción Vite finalizaron correctamente en Docker. \\
  \bottomrule
\end{longtable}

\subsection{20 de septiembre de 2026 - Parametrización y compatibilidad histórica}

\begin{longtable}{@{}p{0.22\textwidth}p{0.18\textwidth}p{0.52\textwidth}@{}}
  \toprule
  \textbf{Actividad} & \textbf{Estado} & \textbf{Resultado} \\
  \midrule
  Catálogo analítico & \statusok & Se separó de Parámetros generales y dispone de permisos propios. Permite CRUD de preguntas por dominio y varias periodicidades soportadas. El objetivo se escribe en cada análisis y las dimensiones quedan bajo propuesta IA, validación determinística y supervisión humana. \\
  Selección de medidas & \statusok & Los identificadores, claves y códigos dejan de competir como candidatos de importe o cantidad; una medida monetaria se elige sólo entre columnas cuantitativas compatibles. \\
  Compatibilidad histórica & \statusok & El campo de dimensiones solicitadas se conserva únicamente para leer contratos anteriores. Las nuevas solicitudes lo envían vacío: el LLM propone las dimensiones desde la necesidad, preguntas y metadatos, y FastAPI sólo valida o completa relaciones técnicas sin imponer el modelo dimensional. \\
  Compatibilidad & \statusok & La verificación distingue errores actuales de diferencias provocadas por versiones anteriores del motor. Una diferencia histórica se informa, pero no retira por sí sola una aprobación válida. \\
  Recuperación & \statusok & Una aprobación retirada puede restaurarse únicamente si supera los controles vigentes, con confirmación de advertencias, justificación y evento de auditoría. \\
  Interfaz & \statusok & Los mensajes indican si debe corregirse, descartarse, crear una versión o confirmar compatibilidad; la evidencia ya no expone referencias al número de un objetivo académico. \\
  Prueba integrada & \statusok & La versión 26 fue reconocida como compatible y su aprobación se restauró con auditoría. Los intentos 29, 30 y 32 se descartaron con motivo trazable. Con el flujo corregido y sin dimensiones solicitadas, Qwen local generó la versión 33 con \file{sales-bi-v3}: propuso por sí mismo producto, cliente y territorio, el hecho de ventas, dos medidas y tres KPI; obtuvo cero errores y advertencias y superó los cuatro controles de evidencia. Gemini volvió a quedar activo. \\
  Pruebas & \statusok & Ruff, Mypy y 48 pruebas Pytest; ESLint, 9 pruebas Vitest y construcción Vite finalizaron correctamente en Docker. \\
  \bottomrule
\end{longtable}

\subsection{21 de septiembre de 2026 - Preparación de KPI sugeridos por IA para Sprint 4}

\begin{longtable}{@{}p{0.22\textwidth}p{0.18\textwidth}p{0.52\textwidth}@{}}
  \toprule
  \textbf{Actividad} & \textbf{Estado} & \textbf{Resultado} \\
  \midrule
  \endfirsthead
  \toprule
  \textbf{Actividad} & \textbf{Estado} & \textbf{Resultado} \\
  \midrule
  \endhead
  KPI dinámicos & \statusok & El LLM propone una cantidad variable de KPI desde la necesidad y metadatos comprobados; no existe un catálogo fijo ni se fuerza una lista en el frontend. \\
  Semántica & \statusok & Cada sugerencia declara medida, dimensiones, período, unidad y receta; clientes nuevos, margen y estados sólo se admiten cuando sus requisitos sean verificables. \\
  IA supervisada & \statusok & FastAPI asignará recetas determinísticas de agregación, razón o participación y no aceptará SQL, fórmulas libres, columnas, relaciones o denominadores inventados. \\
  Experiencia profesional & \statusok & Se especificó un recorrido de selección, revisión de indicadores, confirmación, materialización y validación con estados persistidos, trazabilidad desplegable, recuperación, diseño responsive y accesibilidad. \\
  Validación & \statusok & El expediente de Sprint 4 deberá contrastar cada KPI contra AdventureWorks OLTP y conservar receta, filtros, período, propuesta y ejecución; la demostración final evidenciará al menos cinco KPI sugeridos por IA y conciliados. \\
  Límite & \statusok & No se modificó la entrega ni la aprobación del Sprint 3: el cambio prepara el contrato ejecutable de Sprint 4 y mantiene sus KPI como propuestas estructurales, todavía no calculadas. \\
  \bottomrule
\end{longtable}

\subsection{23 de septiembre de 2026 - Motor ETL, conciliación y proveedor Groq}

\begin{longtable}{@{}p{0.22\textwidth}p{0.18\textwidth}p{0.52\textwidth}@{}}
  \toprule
  \textbf{Actividad} & \textbf{Estado} & \textbf{Resultado} \\
  \midrule
  \endfirsthead
  \toprule
  \textbf{Actividad} & \textbf{Estado} & \textbf{Resultado} \\
  \midrule
  \endhead
  Elegibilidad ETL & \statusok & Las propuestas aprobadas se revalidan contra la instantánea y reglas analíticas vigentes. Se bloquean dimensiones ambiguas, tablas puente, ausencia de atributos descriptivos y contradicciones de granularidad; la interfaz explica cómo corregir cada caso. \\
  Materialización & \statusok & El constructor genera consultas desde columnas y claves foráneas verificadas, carga dimensiones con claves sustitutas, deriva calendario, limpia texto, deduplica y reemplaza transaccionalmente \file{mart\_ventas}; no ejecuta SQL producido por el LLM. \\
  Conciliación & \statusok & La propuesta Groq 48 fue aprobada por la analista y originó la ejecución de referencia 5. Se leyeron y cargaron 121317 filas del hecho, sin rechazos ni diferencias. Las cuatro dimensiones cargaron 504 productos, 19820 clientes, 10 territorios y 1124 fechas distintas. La conciliación respeta la agregación declarada: suma, promedio o conteo distinto. \\
  KPI & \statusok & Groq sugirió cinco KPI materializables y el sistema los compiló a recetas versionadas: ventas totales, unidades vendidas, descuento, precio medio y pedidos distintos. Los resultados fueron 109846381.399888, 274914, 342.8500, 465.0934956741429478 y 31465, respectivamente; todos quedaron conciliados contra el origen. \\
  Español dinámico & \statusok & Esquemas, tablas y columnas mantienen referencia técnica y explicación española. En datos sólo se consideran categorías descriptivas de baja cardinalidad; se excluyen identificadores, PII y texto libre. La IA propone mapeos, pero el analista puede corregirlos, excluirlos y publicarlos sin reemplazar el valor original. Una caída del proveedor permite reintentar esta etapa sin repetir el ETL. \\
  Incidente Gemini & \statusok & Gemini volvió a responder saturado durante la interpretación de la ejecución 3. La carga conciliada se conservó y la plataforma mostró la causa y la acción de recuperación; no se presentó como diferencia numérica. \\
  Groq Cloud & \statusok & Se incorporó, probó y activó el proveedor \file{groq-cloud}, host restringido \file{api.groq.com}, salida JSON Schema estricta, secreto cifrado y niveles de razonamiento compatibles. \file{openai/gpt-oss-120b} generó la propuesta 48 con cero errores y cero advertencias. Se corrigieron límites de salida, semántica de medidas y selección de tablas maestras antes de habilitar la aprobación. \\
  Interfaz & \statusok & El analista dispone de un recorrido guiado para elegir o comparar propuestas, revisar KPI, confirmar transformaciones, ejecutar y comprender la conciliación. Los mensajes diferencian carga fallida, diferencia cuantitativa e interpretación pendiente. La tabla de expedientes recientes permite reabrir una ejecución conservada y continuar su materialización o revisión semántica después de recargar la aplicación o iniciar una nueva sesión, sin repetir el ETL. \\
  Decisión semántica guiada & \statusok & Cada concepto conserva confianza inicial y un expediente determinístico de existencia, columnas, clave y relaciones. Los conceptos de confianza baja o evidencia insuficiente quedan desmarcados; una inclusión excepcional exige confirmación dentro de la plataforma. El analista recibe causa, evidencia y acción sin abrir DBeaver ni escribir SQL. \\
  Copiloto contextual & \statusok & El analista puede consultar a la IA sobre un concepto sin abandonar el asistente. La respuesta persiste conclusión, confianza, evidencia permitida, riesgo, consecuencias y acción recomendada. No recibe filas, credenciales ni SQL libre y nunca altera una selección sin una acción humana explícita. \\
  Pruebas & \statusok & \file{make verify} aprobó Compose y flujo de ramas; Ruff, formato, Mypy y 74 pruebas Pytest; ESLint, 11 pruebas Vitest y construcción Vite; además compiló la documentación en Docker. \\
  \bottomrule
\end{longtable}

\subsubsection{Incidente de migración y resolución}

La primera ejecución aislada de \file{docker compose run --rm migrate} utilizó una imagen local anterior y no encontró la revisión \file{20260918\_07}. No existía un defecto en la cadena Alembic ni pérdida de datos: se confirmó el estado del servicio en ejecución, se aplicó \file{alembic upgrade head} con la imagen vigente del backend y se ejecutó la siembra idempotente. La revisión \file{20260918\_08} quedó aplicada correctamente. En una réplica limpia se debe usar el procedimiento documentado con \file{docker compose up -d --build}, que reconstruye las imágenes antes de migrar.

\subsection{24 de septiembre de 2026 - Copiloto semántico y propuesta autocorregida}

\begin{longtable}{@{}p{0.22\textwidth}p{0.18\textwidth}p{0.52\textwidth}@{}}
  \toprule
  \textbf{Actividad} & \textbf{Estado} & \textbf{Resultado} \\
  \midrule
  \endfirsthead
  \toprule
  \textbf{Actividad} & \textbf{Estado} & \textbf{Resultado} \\
  \midrule
  \endhead
  Copiloto contextual & \statusok & Se incorporó una conversación persistente por concepto con conclusión, confianza, evidencia técnica restringida, riesgo, consecuencias y acción recomendada. Nunca recibe filas, credenciales o SQL libre ni cambia una selección sin acción explícita. \\
  Guardarraíl monetario & \statusok & El proveedor propuso una receta incompleta de tasa por cantidad. El backend ahora reconoce también fórmulas parciales, las reemplaza por precio por tasa por cantidad cuando las referencias son inequívocas y bloquea propuesta y ETL cuando no puede demostrar la base monetaria. \\
  Ciclo de prueba & \statusok & La propuesta 50 fue descartada de manera auditable; la 51 conservó la evidencia de una cuota temporalmente agotada. La propuesta 52 quedó lista para revisión con cero errores y advertencias, seis KPI, Motivo de venta excluido por baja confianza y descuento monetario autocorregido. \\
  Consulta guiada & \statusok & Groq explicó sobre la propuesta 52 por qué conviene excluir Motivo de venta para el objetivo actual, citó únicamente \file{Sales.SalesOrderHeaderSalesReason} y mantuvo \file{selection\_changed=false}. \\
  Ejecución de referencia 6 & \statusok & La propuesta 52 fue aprobada y materializada en un expediente nuevo. Se conciliaron 121317 filas de origen y destino, sin diferencia ni rechazo, cinco tablas y seis KPI. La revisión semántica publicó únicamente las etiquetas españolas de agrupación territorial y conservó sin cambios los códigos de país. \\
  Supervisión guiada & \statusok & Las transformaciones semánticas ahora indican \textbf{Revisar en el paso 5} y explican que no exigen una aprobación en la preparación. Después de cargar, el expediente conduce a Validar resultados e Interpretación semántica para publicar, corregir o excluir cada mapeo. \\
  Prevención de duplicados & \statusok & La propuesta y el conjunto de KPI ya ejecutados muestran el estado y número de expediente, reemplazan el botón de carga por \textbf{Abrir expediente} y no vuelven a materializar el datamart. FastAPI repite el control dentro de una transacción con bloqueo, devuelve una explicación accionable ante solicitudes duplicadas o simultáneas y conserva el reintento sólo para fallos. \\
  Divisa monetaria & \statusok & El sistema dejó de presentar la unidad genérica como resultado final. Comprueba una moneda de origen o base única dentro del alcance relacional, conserva la referencia técnica y formatea los importes con código ISO. En la ejecución 6 verificó \file{USD} mediante \file{Sales.CurrencyRate.FromCurrencyCode} y actualizó los seis KPI sin repetir el ETL; una fuente mixta o ambigua conserva \textit{moneda de origen} y muestra una acción guiada. \\
  Pruebas de regresión & \statusok & Ruff, formato, Mypy y 76 pruebas Pytest; ESLint, 11 pruebas Vitest y construcción Vite aprobaron los cambios de recuperación, claridad, protección operacional y evidencia monetaria. \\
  \bottomrule
\end{longtable}

\subsection{24 de septiembre de 2026 - Normalización documental y Capítulo III}

\begin{longtable}{@{}p{0.22\textwidth}p{0.18\textwidth}p{0.52\textwidth}@{}}
  \toprule
  \textbf{Actividad} & \textbf{Estado} & \textbf{Resultado} \\
  \midrule
  \endfirsthead
  \toprule
  \textbf{Actividad} & \textbf{Estado} & \textbf{Resultado} \\
  \midrule
  \endhead
  Separación editorial & \statusok & El manual, la bitácora y los cinco informes de sprint conservan el detalle técnico completo sin identidad institucional. No se eliminaron decisiones, incidencias, comandos, resultados ni trazabilidad. \\
  Versión académica & \statusok & Se creó un Capítulo III independiente con propuesta tecnológica, factibilidad, metodología, requisitos, casos de uso, UML, arquitectura, sprints, validación, resultados y referencias institucionales. \\
  Ingeniería de software & \statusok & Los diagramas editables cubren ciclo de desarrollo, casos de uso, componentes, despliegue, secuencias y trazabilidad de artefactos. Los casos de uso incluyen actores, precondiciones, excepciones y postcondiciones del ETL. \\
  Automatización documental & \statusok & Make y CI compilan ocho PDF oficiales y verifican que las versiones generadas coincidan con sus fuentes. \\
  Criterio de conservación & \statusok & El capítulo académico resume y enlaza; los informes técnicos permanecen como evidencia integral y anexos del trabajo. \\
  \bottomrule
\end{longtable}

\subsection{25 de septiembre de 2026 - Cierre funcional y documental del Sprint 5}

\begin{longtable}{@{}p{0.22\textwidth}p{0.18\textwidth}p{0.52\textwidth}@{}}
  \toprule
  \textbf{Actividad} & \textbf{Estado} & \textbf{Resultado} \\
  \midrule
  \endfirsthead
  \toprule
  \textbf{Actividad} & \textbf{Estado} & \textbf{Resultado} \\
  \midrule
  \endhead
  Resolución de clientes & \statusok & Se corrigió la dimensión sin sustituirla por todas las personas del origen. El motor conserva a quienes participan en ventas, atraviesa relaciones verificadas y resuelve nombre y tipo de persona u organización según cobertura y calidad descriptiva. \\
  Evidencia ETL & \statusok & La propuesta 54 originó la ejecución 7: 121317 filas de origen y destino, diferencia cero, 19820 clientes con nombre y tipo, 19119 personas y 701 organizaciones. Las versiones 52/54 y ejecuciones 6/7 permanecen trazables. \\
  Analítica por perfil & \statusok & Se implementaron vistas ejecutiva y analítica con filtros de servidor, KPI variables, evolución, productos, clientes, territorios, calidad y hallazgos determinísticos. Los conteos y unidades se etiquetan por significado y los importes usan moneda ISO comprobada. \\
  Copiloto analítico & \statusok & La conversación contextual dispone de guía de preguntas, historial visible, evidencia, cautela y seguimientos. Sólo recibe agregados autorizados; no recibe filas, PII, secretos ni SQL. La auditoría conserva huella y metadatos, no la conversación completa. \\
  Reportería & \statusok & PDF y Excel se reconstruyen en el servidor con la selección y permisos vigentes. El PDF separa seis páginas; el Excel contiene seis hojas con números, moneda, áreas de impresión y gráficos editables. Ambos fueron renderizados para control visual. \\
  Compatibilidad Groq & \statusok & La prueba de \file{openai/gpt-oss-120b} traduce niveles soportados, reserva salida suficiente y usa razonamiento oculto; el error 400 de esfuerzos medio y alto quedó diagnosticado y corregido. \\
  Alcance predictivo & \statusok & El tutor retiró pronóstico, MAPE y RMSE. Se documentó como decisión controlada porque el objetivo es construir y validar un datamart, no entrenar un sistema de predicción. \\
  SDD & \statusok & Se registraron SPR-05-01 a SPR-05-05 para dashboard, reportería, resolución de entidades, compatibilidad LLM y copiloto contextual, con criterios y evidencia verificables. \\
  Documentación & \statusok & Se actualizaron alcance, arquitectura, réplica, validación, README, manual técnico, Capítulo III y bitácora; se creó INF-SPR-05 y se incorporó su PDF a Make y CI. \\
  \bottomrule
\end{longtable}

\section{Matriz de verificación de la fase}

\begin{tabularx}{\textwidth}{@{}X l X@{}}
  \toprule
  \textbf{Control} & \textbf{Estado} & \textbf{Evidencia esperada} \\
  \midrule
  Compose autocontenido válido & \statusok & \file{docker compose config --quiet} \\
  Compose de imágenes válido & \statusok & \file{docker compose -f compose.release.yaml config --quiet} \\
  Backend calidad y pruebas & \statusok & Ruff, formato, Mypy y 84 pruebas Pytest dentro de Docker \\
  Frontend calidad y pruebas & \statusok & ESLint, 11 pruebas Vitest y build de Vite dentro de Docker \\
  PostgreSQL disponible & \statusok & Esquemas \file{app} y \file{mart}; endpoint \file{/health/ready} en \file{ok} \\
  AdventureWorks restaurada & \statusok & 71 tablas; \file{bi\_reader}: lectura 1 e inserción 0 \\
  Interfaz disponible & \statusok & Vite saludable en 5173, Nginx opcional en 8080 y API accesible mediante proxy \\
  Ollama local opcional & \statusok & Perfil \file{local-llm}, modelo ágil \file{qwen2.5:3b} persistente y prueba FastAPI de servicio más modelo, sin puerto público ni descarga CI \\
  PDF legible & \statusok & Todas las páginas renderizadas e inspeccionadas sin cortes ni solapamientos \\
  Git local & \statusok & \file{develop} para integración, \file{main} para entregas y ramas enfocadas por pull request \\
  GitHub/GHCR & \statusok & Repositorio público, dos colaboradoras con escritura, ramas protegidas y cinco imágenes públicas en GHCR; la réplica desde una VM Ubuntu se realizó con Docker y VS Code. \\
  SDD y trazabilidad & \statusok & Especificaciones de Sprints 2 a 5 versionadas; SPR-05-01 a SPR-05-05 cubren dashboard, reportes, calidad semántica, proveedores y conversación. \\
  \bottomrule
\end{tabularx}

\section{Próximos pasos}

\begin{enumerate}
  \item Integrar el Capítulo III en la plantilla definitiva del trabajo de titulación cuando se confirmen tutor, numeración general, estilo de citas y anexos exigidos.
  \item Ejecutar pruebas formales de usabilidad con analistas y usuarios de negocio, y registrar sus instrumentos y resultados.
  \item Realizar juicio de expertos y contraste secundario con AdventureWorksDW únicamente donde exista equivalencia semántica documentada.
  \item Diseñar como extensión futura una actualización periódica de datos, diferenciada de la carga inicial y con controles propios.
  \item Promover a \file{main} únicamente una integración estable, revisada y con CI aprobada.
\end{enumerate}

\section{Historial del documento}

\begin{longtable}{@{}p{0.62in} p{0.86in} p{4.82in}@{}}
  \toprule
  \textbf{Versión} & \textbf{Fecha} & \textbf{Cambio} \\
  \midrule
  \endfirsthead
  \toprule
  \textbf{Versión} & \textbf{Fecha} & \textbf{Cambio} \\
  \midrule
  \endhead
  0.1.0 & 24-08-2026 & Apertura de la bitácora y registro de la fase de ambiente, Docker, Git y DevOps. \\
  0.1.1 & 24-08-2026 & Prueba integral desde volúmenes vacíos, corrección del usuario lector e inicialización del repositorio Git. \\
  0.1.2 & 24-08-2026 & Ejecución desde la ruta definitiva y validación de la sesión GitHub de \file{LFXAS}. \\
  0.1.3 & 24-08-2026 & Repositorio privado, primera CI/CD exitosa, imágenes GHCR y mantenimiento de Actions. \\
  0.1.4 & 24-08-2026 & Informe formal del Sprint 1, manual técnico, capturas e integración de todos los PDF en CI. \\
  0.1.5 & 24-08-2026 & Colaboradores, estrategia \file{develop/main}, protección, control de ramas, Dependabot y reinicio seguro de SQL Server. \\
  0.1.6 & 24-08-2026 & Vista Nginx paralela, entrega GHCR aislada y documentación de ambientes simultáneos con CI/CD. \\
  0.1.7 & 24-08-2026 & Cinco paquetes GHCR públicos, trazabilidad del incidente de autorización y réplica externa pendiente. \\
  0.1.8 & 24-08-2026 & Manual técnico 1.4.0 con operación completa de desarrollo, vista Nginx y entrega GHCR. \\
  0.1.9 & 01-09-2026 & Cierre documental del Sprint 1: manual CI/CD/DevOps ampliado, corrección de visibilidad pública y adopción de SDD con plantilla y especificación del Sprint 2. \\
  0.2.0 & 02-09-2026 & Profundización del manual técnico y del informe del Sprint 1; separación de un manual docente anónimo en LaTeX y registro de los PR de cierre. \\
  0.2.1 & 04-09-2026 & Ampliación SDD del Sprint 2: UI/UX responsive, flujos, navegación, componentes, accesibilidad y criterios verificables previos a la implementación de RBAC. \\
  0.2.2 & 04-09-2026 & Se incorporó la configuración segura y parametrizable de un proveedor LLM para el prototipo, sin adelantar la generación de propuestas BI. \\
  0.2.3 & 04-09-2026 & Se aprobó el catálogo inicial Gemini Cloud, Qwen Cloud y Ollama local; el ADR 0002 define secreto por referencia, adaptadores y \file{qwen3:4b} como opción local. \\
  0.2.4 & 06-09-2026 & Implementación inicial del Sprint 2: migración RBAC/parámetros, JWT, auditoría, adaptadores LLM, cascarón React responsive y verificación integrada en Docker. \\
  0.2.5 & 06-09-2026 & CRUD administrativo del Sprint 2: migración de estado de permisos, edición y desactivación lógica protegidas por RBAC, formularios React y validaciones Docker de backend y frontend. \\
  0.2.6 & 06-09-2026 & Mejora UX del Sprint 2: menús agrupados y administrables por módulo, hamburguesa responsive y paginación servidor-cliente para recursos administrativos. \\
  0.2.7 & 10-09-2026 & Corrección de Sprint 2: protección persistente de administración, selectores guiados, navegación plegable/ocultable, aislamiento de formularios, catálogo de parámetros y dirección visual BI. \\
  0.2.8 & 10-09-2026 & Ajuste UX: listbox de permisos de ancho legible y control lateral compacto superpuesto para ocultar/restaurar la navegación sin desplazarla. \\
  0.2.9 & 10-09-2026 & Corrección UX de creación de usuarios: validaciones API legibles por campo y orientación para el formato de correo. \\
  0.2.10 & 10-09-2026 & Corrección de API: respuesta consistente después de crear o editar usuarios con roles asignados. \\
  0.2.11 & 10-09-2026 & Eliminación física controlada: confirmación de interfaz, bloqueo explícito de dependencias hijas, protección de registros base y auditoría de las eliminaciones permitidas. \\
  0.2.12 & 10-09-2026 & Corrección de asociaciones: permisos de rol mediante \texttt{PUT}, contraseña opcional al editar y tablas de casillas legibles para roles y permisos. \\
  0.2.13 & 10-09-2026 & Corrección UX de LLM: las pruebas de conexión exitosas se comunican con confirmación verde accesible y los fallos conservan alerta roja. \\
  0.2.14 & 11-09-2026 & Alternativa local Ollama: perfil Docker opcional, modelo persistente por equipo, validación de disponibilidad y documentación operativa sin alterar Sprint 1 ni CI/CD. \\
  0.2.15 & 11-09-2026 & Corrección de auditoría para usuarios temporales: preserva eventos y etiqueta histórica del actor, permite eliminación física no protegida mediante \file{ON DELETE SET NULL}. \\
  0.2.16 & 11-09-2026 & Migraciones y catálogo base reproducible: servicio \file{seed} idempotente posterior a Alembic, sin compartir volúmenes ni datos sensibles entre equipos. \\
  0.2.17 & 11-09-2026 & Modelo local predeterminado \file{qwen2.5:3b} y contexto 2048 para una experiencia Ollama ágil; \file{qwen3:4b} queda como alternativa opcional de mayor capacidad. \\
  0.3.0 & 14-09-2026 & Borradores SDD del Sprint 3 para introspección de AdventureWorks, propuesta BI estructurada, validación determinística, aprobación humana y experiencia responsive. \\
  0.3.1 & 14-09-2026 & Redefinición del Sprint 3 para el usuario final: solicitud comercial guiada, alcance técnico automático, glosario español, modo avanzado y futura generación determinística de SQL sin intervención de un programador por análisis. \\
  0.3.2 & 14-09-2026 & Parametrización web de conexiones y secretos, interfaz extensible de conectores e interpretación dinámica de metadatos técnicos a conceptos españoles, sin referencias inventadas ni configuración manual cotidiana. \\
  0.3.3 & 14-09-2026 & Simplificación del Sprint 3 para viabilidad académica: SQL Server único, alcance técnico de sólo lectura, cuatro parámetros reales, credenciales LLM como corrección del Sprint 2 y recorrido centrado en usuario final. \\
  0.3.4 & 18-09-2026 & Especificación UI/UX profesional del Sprint 3: perfiles, wizard inicial, navegación, pantallas, flujo guiado, vista visual del ETL sin drag and drop, responsive, accesibilidad y pruebas de experiencia. \\
  0.3.5 & 18-09-2026 & Cierre pendiente del Sprint 2: credenciales LLM registradas desde la web, cifrado autenticado, clave maestra en volumen separado, migración, pruebas y documentación. \\
  0.3.6 & 18-09-2026 & Inicio funcional del Sprint 3: fuente SQL Server administrada desde la web, prueba de sólo lectura, activación única, secreto cifrado y parámetros tipados. \\
  0.3.7 & 18-09-2026 & Introspección determinística de AdventureWorks: instantánea inmutable, hash estable, deduplicación, búsqueda segura y explorador responsive. \\
  0.3.8 & 19-09-2026 & Cierre funcional del Sprint 3: asistente guiado, interpretación y propuesta IA, validación determinística, revisión humana, trazabilidad y arquitectura extensible. \\
  0.3.9 & 19-09-2026 & Comparación de versiones persistidas, prueba real con Gemini 3.6 Flash y nivel de razonamiento parametrizable por configuración LLM. \\
  0.3.10 & 19-09-2026 & Validación acumulativa visible: integridad, referencias, calidad, coherencia y reejecución determinística; plan de conciliación OLTP--datamart para Sprint 4. \\
  0.3.11 & 19-09-2026 & Corrección de Ollama local: contexto 4096 y mayor presupuesto de salida para evitar el truncamiento del contrato JSON de \file{qwen2.5:3b}. \\
  0.3.12 & 19-09-2026 & Catálogo dinámico de dominios y capacidades, Asistente de datamart genérico, personalización supervisada, historial filtrable y paginado, alcance visible de validación y criterio de selección para Sprint 4. \\
  0.3.13 & 20-09-2026 & Corrección semántica de medidas y KPI, mensajes accionables, reasignación supervisada, retiro de aprobaciones y descarte auditable de versiones. \\
  0.3.14 & 20-09-2026 & Migración correctiva del ciclo de vida de propuestas para persistir los estados de aprobación retirada y descarte sin infringir las restricciones de PostgreSQL. \\
  0.3.15 & 20-09-2026 & Catálogo web de necesidades, selección semántica sin identificadores, compatibilidad histórica y restauración auditada de aprobaciones. \\
  0.3.16 & 20-09-2026 & Corrección determinística de dimensiones solicitadas omitidas por el LLM, depuración auditada de la versión anterior y validación real sin advertencias con Qwen. \\
  0.3.17 & 20-09-2026 & Separación del catálogo analítico por dominio y permisos, CRUD de preguntas, periodicidades múltiples y retiro de dimensiones predefinidas para reforzar la propuesta IA supervisada. \\
  0.3.18 & 20-09-2026 & Validación integrada del flujo sin dimensiones solicitadas: Qwen propuso el modelo de la versión 33, superó los cuatro controles reproducibles y Gemini fue restaurado como proveedor activo. \\
  0.4.0 & 21-09-2026 & Preparación de Sprint 4: contrato de KPI variables sugeridos por IA, recetas determinísticas de agregación, razón y participación, límite técnico ampliado y experiencia profesional de materialización y validación. \\
  0.4.1 & 23-09-2026 & Sprint 4 funcional: selección compatible, materialización transaccional, KPI, conciliación real, interpretación española supervisada y proveedor Groq GPT-OSS 120B. \\
  0.4.2 & 23-09-2026 & Propuesta Groq 48 aprobada y ejecución ETL 5 conciliada; métricas por agregación, dimensiones maestras y cinco KPI verificados. \\
  0.4.3 & 23-09-2026 & Recuperación de expedientes ETL desde el historial: acceso directo a ejecuciones preparadas, conciliadas o pendientes de interpretación sin volver a ejecutar la carga. \\
  0.4.4 & 23-09-2026 & Corrección de persistencia de la revisión semántica: las etiquetas publicadas conservan estado, responsable y comentario; el expediente validado muestra un comprobante final y no reinicia la confirmación. \\
  0.4.5 & 23-09-2026 & Control semántico de medidas calculadas: detección de tasas usadas como importes, autocorrección declarativa, operaciones aritméticas seguras y propuesta 49 derivada de la 48 con descuento monetario reproducible. \\
  0.4.6 & 23-09-2026 & Validación semántica guiada: evidencia estructural por concepto, exclusión preventiva de confianza baja o insuficiente y confirmación explícita para inclusiones excepcionales sin SQL externo. \\
  0.4.7 & 23-09-2026 & Copiloto contextual auditable para resolver decisiones semánticas con evidencia restringida, riesgos y consecuencias, sin filas, credenciales, SQL libre ni cambios silenciosos. \\
  0.4.8 & 24-09-2026 & Propuesta 52 generada con Groq: seis KPI, concepto de baja confianza excluido, consulta contextual persistida y autocorrección determinística de una receta monetaria parcial. \\
  0.4.9 & 24-09-2026 & Ejecución ETL 6 conciliada y revisión semántica publicada; supervisión ubicada explícitamente en el paso 5 y protección integral contra materializaciones idénticas desde la interfaz y la API. \\
  0.4.10 & 24-09-2026 & Comprobación estándar de divisa monetaria desde evidencia relacional, formato ISO trazable y enriquecimiento de la ejecución 6 como USD sin repetir el ETL. \\
  0.4.11 & 24-09-2026 & Separación entre documentación técnica completa y Capítulo III académico; normalización profesional, UML, casos de uso y compilación de siete PDF en CI. \\
  0.5.0 & 25-09-2026 & Sprint 5: resolución descriptiva de clientes, propuesta 54, ejecución 7, paneles por perfil, copiloto analítico, PDF/Excel fieles, cinco especificaciones, INF-SPR-05 y decisión controlada de retirar pronóstico. \\
  \bottomrule
\end{longtable}

\end{document}
